\part{\centering Weave diagrammatics, braid varieties, and cluster structures}

The object of this first part of the notes is to present the diagrammatic calculus of Demazure weaves and develop its applications to the study of braid varieties and their cluster structures. Specifically, a central result of this first part is the construction of a cluster algebra structure on the coordinate ring of a braid variety, cf.~\Cref{table:PartI_introduction}. The construction we present builds on the study of Demazure weaves, which can be depicted as certain diagrams on a surface, see e.g.~\Cref{fig:Example_TwoWeavesSL3SL4}. As we will explain, different moves, operations and combinatorial constructions on these diagrams lead to a fruitful interaction between combinatorics, algebra, and geometry. In particular, the relation between Demazure weaves and cluster structures leads to new results, both in the context of cluster algebras and in areas where braid varieties appear for different reasons, such as representation theory or contact and symplectic topology.